\chapter{History of Poland}

\section{Minimal primer}

The conventional beginning of the Polish state is the rule of Mieszko I of the Piast dynasty. His baptism in 966 is traditionally taken as a foundational date, tying his polity to Latin Christianity. Poland became a kingdom in 1025.

\subsection{Poland--Lithuania}

From the late fourteenth century Poland was joined dynastically to the Grand Duchy of Lithuania. Lithuania then included not only modern Lithuania but large parts of present-day Belarus and Ukraine.

For much of the fifteenth and sixteenth centuries the union was ruled by the Jagiellonian dynasty. Around the time of Machiavelli (1469--1527), the Kingdom of Poland and the Grand Duchy of Lithuania remained distinct states sharing a monarch. Members of the Jagiellonian dynasty also ruled Bohemia and Hungary.

The Union of Lublin (1569) created the Polish--Lithuanian Commonwealth. In the sixteenth and seventeenth centuries it was one of Europe's largest states. Its nobility (\emph{szlachta}) enjoyed extensive political privileges and royal power was substantially limited; its parliament was the \emph{Sejm}.

\subsection{Partitions}

During the eighteenth century the Commonwealth weakened while its neighbors grew stronger. Russia, Prussia and Austria partitioned it in 1772, 1793 and 1795. After the third partition Poland ceased to exist as an independent state.

Polish lands remained divided among neighboring empires for 123 years, although Polish national identity survived and several uprisings sought to restore independence.

\subsection{Independence and the Second World War}

Poland regained independence in 1918 after the First World War and subsequently fought Soviet Russia in the Polish--Soviet War (1919--1921).

Nazi Germany invaded Poland from the west on 1 September 1939. The Soviet Union invaded from the east on 17 September, following the secret territorial arrangements of the Molotov--Ribbentrop Pact. Nazi occupation brought mass murder and the Holocaust; Soviet repression included the Katyn massacre of roughly 22,000 Polish officers, police and members of the intelligentsia in 1940.

\subsection{Communist Poland and the present}

After the Second World War Poland became a communist state aligned with and heavily constrained by the Soviet Union, and later belonged to the Warsaw Pact. The independent trade-union movement Solidarność (Solidarity), founded in 1980 and associated especially with Lech Wałęsa, became a mass opposition movement. Negotiations and partially free elections in 1989 led to the end of communist rule.

Poland joined NATO in 1999 and the European Union in 2004.

A useful geographical intuition is that Poland lies on the North European Plain, between the German lands to the west and Russia and eastern Europe to the east, with relatively few major natural barriers. Its territory has repeatedly been crossed or contested by powerful neighbors. But Poland was not merely a disputed territory: Poland--Lithuania was itself a major European power in the early modern period.

The partitions, the Soviet invasion of 1939 and postwar Soviet domination are essential background for contemporary Polish security policy toward Russia. Russia's full-scale invasion of Ukraine in 2022 therefore had particular resonance in Poland, which became one of Ukraine's major European supporters and regards NATO membership as a central security guarantee.
